\documentclass[11pt]{article}
\usepackage{amsmath,amssymb,booktabs,geometry,siunitx}
\usepackage[hidelinks]{hyperref}
\geometry{margin=1in}
\begin{document}

\title{Numerical methods and discrete reference oracles}
\author{cse-heat-equation-example}
\date{}
\maketitle

\section{Uniform grid}

Both methods use
\[
  \xi_i=i\Delta\xi,\qquad Fo^n=n\Delta Fo,
  \qquad i=0,\ldots,N,
\]
with \(\Delta\xi=1/N\), an integer number of equal time steps, exact zero
boundaries, and
\[
  r=\frac{\Delta Fo}{(\Delta\xi)^2}.
\]

\section{M-FTCS-001: forward-time centered-space}

The project-derived explicit update is
\begin{equation}
  \theta_i^{n+1}
  =\theta_i^n+r\left(
    \theta_{i-1}^n-2\theta_i^n+\theta_{i+1}^n
  \right),
  \qquad i=1,\ldots,N-1.
  \label{eq:ftcs}
\end{equation}
It is first order in \(\Delta Fo\), second order in \(\Delta\xi\), and
linearly stable for this diffusion stencil only when
\begin{equation}
  r\leq\frac{1}{2}.
  \label{eq:ftcs-stability}
\end{equation}
\texttt{solve\_ftcs} rejects a larger value before time integration.

\section{M-CN-001: Crank-Nicolson}

Trapezoidal centering of the same spatial difference gives
\begin{equation}
\begin{split}
  \theta_i^{n+1}-\theta_i^n
  =\frac{r}{2}\big[&
    \theta_{i-1}^{n+1}-2\theta_i^{n+1}+\theta_{i+1}^{n+1}\\
    &+\theta_{i-1}^{n}-2\theta_i^n+\theta_{i+1}^n
  \big].
\end{split}
\label{eq:crank-nicolson}
\end{equation}
The historical finite-difference form appears in Crank and Nicolson (1947),
Section 4, equation (16), printed pages 54--55, DOI
\texttt{10.1017/S0305004100023197}. Zotero parent \texttt{NPC9836V}, PDF
child \texttt{R4GC8HTF}, and Markdown child \texttt{I5DDQJH8} identify the
exact evidence chain.

The implementation solves the tridiagonal interior system with
\texttt{scipy.linalg.solve\_banded}. It records the maximum algebraic residual
over all time steps. The method is second order in both increments for the
smooth declared case. Its linear \(L^2\) stability for positive \(r\) does not
imply unconditional monotonicity. Crank and Nicolson discuss the damping of
Method III error modes and slowly disappearing oscillation for large step
ratio in Section 7(b), printed page 59.

\section{Discrete amplification oracles}

For phase \(\mu\), the centered difference contributes
\(-4\sin^2(\mu/2)\). Independent substitution into the two declared schemes
gives
\begin{align}
  G_{\mathrm{FTCS}}(\mu,r)
  &=1-4r\sin^2(\mu/2),
  \label{eq:ftcs-amplification}\\
  G_{\mathrm{CN}}(\mu,r)
  &=\frac{1-2r\sin^2(\mu/2)}
          {1+2r\sin^2(\mu/2)}.
  \label{eq:cn-amplification}
\end{align}
These oracles are evaluated separately from the stencil and linear-solve code.
For the worst grid mode, \(\mu=\pi\), FTCS reaches \(-1\) at \(r=1/2\) and
then becomes unstable. Crank-Nicolson remains within unit magnitude for every
positive \(r\), but its factor becomes negative above \(r=1/2\), so a
high-frequency component can alternate in sign.

\section{Implementation and tests}

The public methods and oracles live in
\texttt{src/heat\_verification/model.py}. The compact test matrix in
\path{tests/test_scientific_behavior.py} protects the reference case,
discrete-mode reproduction, exact boundaries, FTCS rejection, the stability
versus monotonicity distinction, and entry into the second-order combined
refinement regime.

\end{document}
